\section{Pulsar: osservazioni, dinamica e struttura stellare}

\subsection{6.1 -- Scoperta e interpretazione}

\begin{frame}[t]{Capitolo 6: obiettivi e convenzioni}
\small
\begin{itemize}
\item Sequenza del libro: scoperta, rotazione e timing, magnetosfera, mezzi interstellari, masse binarie e test della relativita generale.
\item Approfondimenti degli appunti: rette nel piano $P$--$\dot P$, \emph{death line}, dipendenza dal modello del campo magnetico, momento d'inerzia in RG e vincoli sulla EOS.
\item Distinguere $P$ (periodo di rotazione), $\mathcal P$ (potenza emessa), $P_b$ (periodo orbitale); $\nu=1/P$ e $\Omega=2\pi\nu=2\pi/P$.
\item Formule in cgs/Gauss quando compare il campo $B$; altrimenti sono esplicitate $G,c$ e le unita naturali.
\end{itemize}
\medskip
\textbf{Nota:} i dati di pulsar e le figure originali riprendono la \emph{versione storica} del capitolo 6, non un catalogo aggiornato.
\end{frame}

\begin{frame}[t]{6.1 -- Dal segnale periodico alla pulsar}
\small
Nel 1967 Jocelyn Bell Burnell identifico segnali radio di periodicita eccezionale durante studi delle scintillazioni interplanetarie.
\begin{itemize}
\item L'individuazione di sorgenti diverse escludeva un'origine terrestre o artificiale unica.
\item La regolarita rivelo un sistema astronomico con una scala dinamica piccola e un orologio stabile.
\item La spiegazione successiva e una stella di neutroni magnetizzata, in rapida rotazione: un fascio di emissione attraversa periodicamente la linea di vista.
\end{itemize}
\[
\text{rotazione della stella}\quad\longrightarrow\quad\text{modulazione del fascio radio}.
\]
Il periodo del segnale coincide normalmente con il periodo di rotazione (a parte geometrie con piu picchi per giro).
\end{frame}

\subsection{6.2 -- Perche sono stelle di neutroni rotanti}

\begin{frame}[t]{Vincolo di causalita sulla regione emittente}
\small
Se un impulso dura $\Delta t$, una regione che varia in modo causalmente connesso ha dimensione indicativa
\[
\ell\lesssim c\Delta t.
\]
Per la scala del libro, $\Delta t\simeq0.016\,$s:
\[
\ell\lesssim(3\times10^5\,\mathrm{km/s})(0.016\,\mathrm s)
\simeq4.8\times10^3\,\mathrm{km}.
\]
\begin{itemize}
\item Il risultato riguarda la dimensione della \emph{regione variabile}, non necessariamente l'intero corpo celeste.
\item Le dimensioni di una stella di neutroni ($\sim10$--$15\,$km) sono ampiamente compatibili.
\end{itemize}
\end{frame}

\begin{frame}[t]{Tempo dinamico e oscillazioni}
\small
Per una sfera di massa $M$ e raggio $R$, il bilancio idrostatico d'ordine di grandezza e
\[
P_c\sim\frac{GM^2}{R^4},\qquad \bar\rho\sim\frac{M}{R^3}.
\]
Una stima della velocita del suono e $c_s^2\sim P_c/\bar\rho$, da cui
\[
c_s^2\sim\frac{GM}{R},\qquad
\tau_{\rm osc}\sim\frac R{c_s}
\sim\sqrt{\frac{R^3}{GM}}
\sim\frac1{\sqrt{G\bar\rho}}.
\]
\begin{itemize}
\item La scala dipende dalla densita media; i modi normali effettivi introducono coefficienti dipendenti dalla struttura.
\item Una singola frequenza caratteristica di oscillazione non spiega naturalmente la varieta dei periodi delle pulsar.
\end{itemize}
\end{frame}


\begin{frame}[t]{Haensel \S6.12: equilibrio di una stella rotante}
\small
La rotazione rigida stazionaria in relativita' generale richiede una metrica assialsimmetrica, non la soluzione TOV sferica.
Per un fluido barotropico, la pseudoentalpia
\[
H(P)=\int_0^P\frac{dP'}{\varepsilon(P')+P'}
\]
permette di scrivere un primo integrale relativistico dell'equilibrio:
\[
H+\Phi-\ln\Gamma=\text{costante},
\qquad\Gamma=(1-U^2/c^2)^{-1/2}.
\]
$\Phi$ e $U$ dipendono dalla metrica e dalla posizione nella stella: occorre risolvere un problema bidimensionale.
\par\medskip
\footnotesize Haensel et al. (2006), eqs.~(6.95)--(6.102).
\end{frame}

\begin{frame}[t]{Haensel \S6.12: stabilita' e mass shedding}
\begin{columns}[T,onlytextwidth]
\column{0.54\textwidth}
\begin{figure}[!ht]\includegraphics[trim={0cm 0cm 0 0},clip,keepaspectratio,width=0.95\textwidth]{haensel07/M_NS-R_eq}\label{fig:haensel07_M_NS-R_eq}\end{figure}
\column{0.44\textwidth}
\small
Le configurazioni stabili sono limitate da due meccanismi diversi:
\begin{itemize}
\item instabilita' assialsimmetrica a densita' centrale elevata;
\item perdita di massa all'equatore al limite di Keplero.
\end{itemize}
Le sequenze a momento angolare costante sono essenziali per discutere la stabilita'.
\end{columns}
\footnotesize Haensel et al. (2006), Fig.~6.18, \S6.12.2; EOS SLy.
\end{frame}

\begin{frame}[t]{Haensel \S6.12: rotazione e deformazione centrifuga}
\begin{columns}[T,onlytextwidth]
\column{0.53\textwidth}
\begin{figure}[!ht]\includegraphics[trim={0cm 0cm 0 0},clip,keepaspectratio,width=0.95\textwidth]{haensel07/NS_DeltaR-oblat_fHz}\label{fig:haensel07/NS_DeltaR-oblat_fHz}\end{figure}
\column{0.45\textwidth}
\small
Al crescere della frequenza:
\begin{itemize}
\item il raggio equatoriale aumenta;
\item il raggio polare diminuisce;
\item la densita' centrale tende a diminuire a massa barionica fissata.
\end{itemize}
La deformazione e' $O(\Omega^2)$ nel limite lento, ma diventa grande vicino a $\Omega_K$.
\end{columns}
\footnotesize Haensel et al. (2006), Fig.~6.19, \S6.12.2.
\end{frame}

\begin{frame}[t]{Haensel \S6.12: massa massima e rotazione differenziale}
\small
La stima newtoniana del limite centrifugo e'
\[
\Omega_K\sim\sqrt{GM/R^3},\qquad P_K\sim2\pi\sqrt{R^3/(GM)}.
\]
Per EOS barioniche realistiche considerate nel testo, una relazione empirica per rotazione rigida e'
\[
M_{\max}^{\rm rot}\simeq1.18\,M_{\max}^{\rm stat},
\]
che \emph{non} si applica universalmente alle stelle autolegate di materia strana.
La rotazione differenziale puo' sostenere temporaneamente masse maggiori, ma il trasporto del momento angolare ne modifica la stabilita'.
\par\medskip
\footnotesize Haensel et al. (2006), \S6.12.4--6.12.5, eq.~(6.107). Coefficiente riferito ai modelli del 2006.
\end{frame}

\begin{frame}[t]{Limite di rotazione (mass shedding)}
\small
Bilanciando gravitazione e accelerazione centrifuga all'equatore:
\[
\frac{GMm}{R^2}\sim m\Omega_K^2R
\quad\Longrightarrow\quad
\Omega_K\sim\sqrt{\frac{GM}{R^3}}.
\]
Segue un periodo minimo stimato
\[
P_{\min}\sim2\pi\sqrt{\frac{R^3}{GM}}.
\]
Con $M=1.4\,M_\odot$, $R=10\,$km, la stima newtoniana da circa $0.46\,$ms.
\begin{itemize}
\item Un calcolo relativistico di una stella rapidamente rotante modifica il coefficiente (deformazione, \emph{frame dragging}, EOS).
\item Nane bianche di raggio terrestre non possono raggiungere normalmente periodi di rotazione di pochi millisecondi.
\end{itemize}
\end{frame}

\begin{frame}[t]{Modello del faro magnetico}
\small
\begin{columns}[T,onlytextwidth]
\column{.49\textwidth}
\textbf{Geometria}
\begin{itemize}
\item L'asse magnetico forma un angolo $\alpha$ con l'asse di rotazione.
\item La regione emittente e connessa a linee di campo magnetico.
\item L'osservatore registra un impulso quando il fascio attraversa la propria linea di vista.
\end{itemize}
\column{.49\textwidth}
\textbf{Conseguenze}
\begin{itemize}
\item $P$ e molto stabile, ma aumenta lentamente per perdita di momento angolare.
\item La periodicita e un orologio di rotazione: pulsar isolate e binarie diventano sonde di gravita e struttura della materia.
\end{itemize}
\end{columns}
\[
\Omega(t)=\frac{2\pi}{P(t)},\qquad\dot P>0\ \Longleftrightarrow\ \dot\Omega<0.
\]
\end{frame}

\subsection{6.3 e 6.4 -- Timing e famiglie osservative}

\begin{frame}[t]{Osservabili fondamentali del timing}
\small
\begin{itemize}
\item Periodo $P$ e derivata $\dot P$; in alternativa frequenza $\nu$ e derivata $\dot\nu$.
\item Fase rotazionale $\phi(t)=\phi_0+\Omega_0(t-t_0)+\tfrac12\dot\Omega_0(t-t_0)^2+\cdots$.
\item I tempi di arrivo degli impulsi (TOA) vanno corretti al baricentro del Sistema solare, alla dispersione radio e alla dinamica orbitale.
\item In un sistema binario l'effetto Doppler e le accelerazioni apparenti possono alterare $\dot P$ osservata.
\end{itemize}
\[
\nu=\frac1P,\qquad\dot\nu=-\frac{\dot P}{P^2},\qquad
\dot\Omega=-2\pi\frac{\dot P}{P^2}.
\]
\end{frame}

\begin{frame}[t]{Nomi e distribuzione delle pulsar}
\small
La sigla \texttt{PSR} e seguita dalle coordinate equatoriali della sorgente; le designazioni $B$ e $J$ fanno riferimento a epoche coordinate B1950 e J2000.
\begin{itemize}
\item Pulsar giovani: periodo di decine--centinaia di millisecondi, $\dot P$ relativamente grande.
\item Pulsar riciclate: rotazione accelerata da accrescimento in un sistema binario; periodi anche di pochi millisecondi.
\item Sistemi binari relativistici: offrono massa e parametri post-kepleriani.
\item Gli oggetti citati nel libro sono esempi \emph{al momento della sua stesura}; i cataloghi si aggiornano continuamente.
\end{itemize}
\end{frame}

\begin{frame}[t]{6.4 -- Uno zoo di stelle di neutroni}
\small
\begin{columns}[T,onlytextwidth]
\column{.5\textwidth}
\begin{itemize}
\item \textbf{Radio pulsar:} emissione pulsata alimentata prevalentemente dalla rotazione.
\item \textbf{Magnetar:} emissione X/soft gamma spesso alimentata dall'energia magnetica; $B$ caratteristico molto elevato.
\item \textbf{CCO:} sorgenti termiche nei resti di supernova.
\end{itemize}
\column{.5\textwidth}
\begin{itemize}
\item \textbf{RRAT:} impulsi radio sporadici.
\item \textbf{INS termiche:} neutron stars isolate, tipicamente poco luminose in radio.
\item \textbf{Binarie X:} accrescimento da una compagna; talvolta pulsazioni e burst.
\end{itemize}
\end{columns}
\medskip
Le classi sono \emph{osservative}: la stessa stella di neutroni puo mostrare comportamenti differenti a seconda di campo, temperatura, accrescimento e geometria.
\end{frame}

\section{Rotazione, momento di inerzia ed emissione}

\subsection{6.5.1 -- Energia di rotazione e momento d'inerzia}

\begin{frame}[t]{Energia di rotazione: variabili del modello}
\small
Per una stella che ruota lentamente e in modo quasi rigido:
\[
E_{\rm rot}=\frac12 I\Omega^2
=\frac{2\pi^2 I}{P^2},\qquad J=I\Omega.
\]
\begin{itemize}
\item $I$ e il momento d'inerzia rispetto all'asse di rotazione.
\item L'energia radiata, il vento di particelle e altre coppie esterne riducono $J$ e $E_{\rm rot}$.
\item La dipendenza $I=I(M,R,\mathrm{EOS},\Omega)$ va tenuta distinta dall'approssimazione $I=$ costante.
\end{itemize}
Se $I$ e costante,
\[
\dot E_{\rm rot}=I\Omega\dot\Omega
=-\frac{4\pi^2 I\dot P}{P^3}<0.
\]
\end{frame}

\begin{frame}[t]{Momento d'inerzia: derivazione in meccanica classica}
\small
Per una distribuzione sferica $\rho(r)$ e rotazione attorno a $z$:
\[
I_z=\int \rho(r)r_\perp^2\,dV,
\qquad r_\perp=r\sin\theta.
\]
Con $dV=r^2\sin\theta\,dr\,d\theta\,d\varphi$:
\[
\begin{aligned}
I_z&=2\pi\int_0^R\rho(r)r^4dr
\int_0^\pi\sin^3\theta\,d\theta\\
&=\frac{8\pi}{3}\int_0^R\rho(r)r^4dr.
\end{aligned}
\]
Il peso $r^4$ rende $I$ particolarmente sensibile alla distribuzione radiale della materia.
\end{frame}

\begin{frame}[t]{Sfera omogenea: ricavare $I=2MR^2/5$}
\small
Per $\rho(r)=\rho_0$:
\[
I=\frac{8\pi}{3}\rho_0\int_0^Rr^4dr
=\frac{8\pi}{15}\rho_0R^5.
\]
Poiche $M=(4\pi/3)\rho_0R^3$:
\[
I=\frac25 MR^2\qquad\text{(sfera omogenea, Newton)}.
\]
Con $M=1.4M_\odot$, $R=10\,$km,
\[
I\simeq1.11\times10^{38}\,\mathrm{kg\,m^2}
\simeq1.11\times10^{45}\,\mathrm{g\,cm^2}.
\]
Un profilo centralmente condensato e gli effetti relativistici modificano il fattore $I/(MR^2)$.
\end{frame}

\begin{frame}[t]{Perche $I$ dipende dalla EOS}
\small
\begin{itemize}
\item La soluzione TOV determina la funzione $\varepsilon(r)$ e il raggio $R$ per una massa gravitazionale $M$.
\item EOS diverse distribuiscono diversamente l'energia e la pressione in funzione di $r$.
\item A parita di $M$, raggi e profili differenti producono momenti d'inerzia differenti.
\end{itemize}
\[
I\sim k_I MR^2,\qquad k_I\ \text{non e universale}.
\]
La misura di $I$ in un sistema binario relativistico, insieme con $M$, offre percio un vincolo aggiuntivo su $R$ e sull'EOS.
\medskip
\textbf{Attenzione:} nelle stelle di neutroni l'integrale newtoniano non e quantitativamente sufficiente.
\end{frame}

\begin{frame}[t]{Momento d'inerzia in RG: definizione}
\small
In un campo gravitazionale forte, per \emph{rotazione lenta}:
\[
I\equiv\left(\frac{\partial J}{\partial\Omega}\right)_{\Omega\to0}
\qquad\Rightarrow\qquad J=I\Omega+O(\Omega^3).
\]
La metrica statica di partenza e
\[
ds^2=e^{\nu(r)}c^2dt^2-e^{\lambda(r)}dr^2-r^2d\Omega_2^2,
\]
con $e^{-\lambda}=1-2Gm(r)/(rc^2)$ (firma $+---$).
\begin{itemize}
\item La rotazione aggiunge $g_{t\varphi}$ e una velocita angolare locale di trascinamento $\omega(r)$.
\item La materia si muove rispetto ai riferimenti inerziali locali con $\bar\omega=\Omega-\omega$.
\end{itemize}
\end{frame}

\begin{frame}[t]{Hartle: equazione di frame dragging}
\small
In unita $G=c=1$, al primo ordine in $\Omega$ definire
\[
j(r)=e^{-[\nu(r)+\lambda(r)]/2},\qquad
\bar\omega(r)=\Omega-\omega(r).
\]
L'equazione per la rotazione lenta e
\[
\frac1{r^4}\frac{d}{dr}
\left(r^4 j\frac{d\bar\omega}{dr}\right)
+\frac4r\frac{dj}{dr}\bar\omega=0.
\]
\begin{itemize}
\item Regolarita al centro: $d\bar\omega/dr=0$ in $r=0$.
\item Nel vuoto esterno: $\bar\omega=\Omega-2J/r^3$.
\item Il confronto fra soluzione interna ed esterna determina $J$ e quindi $I=J/\Omega$.
\end{itemize}
\end{frame}

\begin{frame}[t]{Hartle: integrale relativistico per $I$}
\small
Per fluido perfetto lentamente rotante, in unita $G=c=1$:
\[
I=\frac{8\pi}{3}\int_0^Rdr\,r^4(\varepsilon+P)
 e^{(\lambda-\nu)/2}\frac{\bar\omega(r)}{\Omega}.
\]
Qui $\varepsilon,P$ sono nella soluzione TOV, $\nu,\lambda$ le funzioni metriche e $\bar\omega$ viene dall'equazione del trascinamento.
\medskip
\textbf{Limite debole:} $P\ll\varepsilon$, $\nu,\lambda,\omega\to0$ e $\varepsilon\to\rho$ in $c=1$:
\[
I\longrightarrow\frac{8\pi}{3}\int_0^R\rho(r)r^4dr.
\]
Questo rende esplicito il collegamento fra il capitolo 2 (RG/TOV) e il timing delle pulsar.
\end{frame}


\begin{frame}[t]{Haensel \S6.10: rotazione lenta e ordine perturbativo}
\small
Il parametro che controlla la deformazione centrifuga e'
\[
\epsilon_{\rm rot}=\frac{R^3\Omega^2}{GM}\ll1.
\]
L'inversione del verso di rotazione lascia invariata la struttura scalare:
\[
M(\Omega)=M(0)+O(\Omega^2),\quad R(\Omega)=R(0)+O(\Omega^2).
\]
Invece il momento angolare cambia segno: $J=I\Omega+O(\Omega^3)$.
Quindi $I=J/\Omega$ si calcola, all'ordine dominante, sul fondo sferico TOV.
\par\medskip
\footnotesize Fonte: Haensel, Potekhin e Yakovlev (2006), \S6.10, pp.~323--324.
\end{frame}

\begin{frame}[t]{Haensel \S6.10: matching del frame dragging}
\small
Indichiamo con $\varpi(r)=\Omega-\omega_{\rm FD}(r)$ la velocita' angolare rispetto ai riferimenti inerziali locali.
Nel vuoto esterno, all'ordine lineare in $\Omega$,
\[
\omega_{\rm FD}(r)=\frac{2GJ}{c^2r^3},\qquad
\varpi(r)=\Omega-\frac{2GJ}{c^2r^3}.
\]
La continuita' alla superficie implica
\[
\Omega=\varpi(R)+\frac{2GJ}{c^2R^3},\qquad I=\frac{J}{\Omega}.
\]
La soluzione interna richiede le equazioni TOV e l'equazione differenziale di Hartle; non basta la distribuzione newtoniana della massa.
\par\medskip
\footnotesize Fonte: Haensel et al., eqs.~(6.65)--(6.70).
\end{frame}

\begin{frame}[t]{Haensel \S6.10: $I(M)$ e dipendenza dall'EOS}
\begin{columns}[T,onlytextwidth]
\column{0.55\textwidth}
\begin{figure}[!ht]\includegraphics[trim={0cm 0cm 0 0},clip,keepaspectratio,width=0.95\textwidth]{haensel07/NS-momentinertia}\label{fig:haensel07_NS-momentinertia}\end{figure}
\column{0.43\textwidth}
\small
La figura confronta il momento d'inerzia di stelle lentamente rotanti per diverse EOS.
\begin{itemize}
\item A massa fissata, EOS diverse implicano raggi e profili di densita' diversi.
\item Il massimo di $I(M)$ non coincide in generale con quello di $M$.
\item L'osservazione di $I$ restringerebbe i modelli della materia densa.
\end{itemize}
\end{columns}
\footnotesize Haensel et al. (2006), Fig.~6.16, p.~325 (modelli e dati storici).
\end{frame}

\begin{frame}[t]{Haensel \S6.10: compattezza e $I/(MR^2)$}
\begin{columns}[T,onlytextwidth]
\column{0.52\textwidth}
\begin{figure}[!ht]\includegraphics[trim={0cm 0cm 0 0},clip,keepaspectratio,width=0.95\textwidth]{haensel07/I_MR2-rg_R}\label{fig:haensel07_I_MR2-rg_R}\end{figure}
\column{0.46\textwidth}
\small
Definiamo $x_{\rm GR}=2GM/(Rc^2)$. La figura mostra una correlazione approssimativamente universale fra
\[
\frac{I}{MR^2}\quad\text{e}\quad x_{\rm GR}.
\]
La dispersione residua dipende dall'EOS. La relazione e' una calibrazione empirica, non un'identita' esatta della relativita' generale.
\end{columns}
\footnotesize Haensel et al. (2006), Fig.~6.17 e \S6.10.4, pp.~326--327.
\end{frame}

\begin{frame}[t]{Haensel \S6.10: quanto conta la crosta?}
\small
Il contributo della crosta e' la porzione dell'integrale relativistico compresa fra $R_{\rm crust}$ e $R$:
\[
I_{\rm crust}=\frac{8\pi}{3}\int_{R_{\rm crust}}^R
r^4\left(\rho+\frac{P}{c^2}\right)
\frac{\varpi(r)}{\Omega}e^{-\Phi-\lambda}\,dr.
\]
Per una crosta sottile, $P\ll\rho c^2$, una prima stima e'
\[
I_{\rm crust}\simeq\frac23 M_{\rm crust}R^2\frac{\varpi(R)}{\Omega}.
\]
L'approssimazione mette in evidenza il fattore $r^2$: una massa crostale piccola puo' contribuire significativamente a $I$.
\par\medskip
\footnotesize Haensel et al. (2006), eqs.~(6.71)--(6.73); attenzione alle convenzioni metriche.
\end{frame}

\begin{frame}[t]{Spin-down luminosity e caso Crab}
\small
Definire la potenza liberata dalla frenata
\[
\dot E_{\rm sd}\equiv-\dot E_{\rm rot}
=4\pi^2I\frac{\dot P}{P^3}.
\]
Con $P=0.0334\,$s, $\dot P=4.21\times10^{-13}$ e $I=10^{45}\,\mathrm{g\,cm^2}$:
\[
\dot E_{\rm sd}\approx4.5\times10^{38}\,\mathrm{erg\,s^{-1}}.
\]
L'energia va in emissione magnetosferica, vento relativistico e nebulosa; non coincide in generale con la luminosita radio.
\begin{itemize}
\item Misurare $P,\dot P$ fornisce una scala di potenza \emph{assumendo} $I$.
\item Un cambiamento temporale di $I$ introdurrebbe un termine addizionale $\tfrac12\dot I\Omega^2$ in $\dot E_{\rm rot}$.
\end{itemize}
\end{frame}


\begin{frame}[t]{Haensel \S9.5: bilancio energetico della nebulosa del Crab}
\small
L'energia rotazionale persa dalla pulsar alimenta, fra altri canali, la nebulosa:
\[
-\dot E_{\rm rot}=I\Omega|\dot\Omega|
=\frac{4\pi^2I\dot P}{P^3}.
\]
Un bilancio schematico della potenza richiesta e'
\[
I\Omega|\dot\Omega|\ \gtrsim\ L_{\rm rad}+\dot E_{\rm kin}+\dot E_{\rm altre},
\]
con termini di accelerazione dei filamenti, radiazione e altre perdite da modellare.
\begin{itemize}
\item Il timing determina $\Omega$ e $\dot\Omega$.
\item La stima energetica della nebulosa vincola $I$.
\item Il risultato dipende da distanza, massa e dinamica della nebulosa.
\end{itemize}
\footnotesize Haensel et al. (2006), \S9.5.1--9.5.3, pp.~501--504; stime storiche.
\end{frame}

\subsection{6.5.2 -- Dipolo magnetico e campo caratteristico}

\begin{frame}[t]{Campo di un dipolo magnetico}
\small
In unita gaussiane, per un momento magnetico $\boldsymbol\mu$ nel vuoto:
\[
\mathbf B(\mathbf r)=\frac{3\hat{\mathbf r}
(\boldsymbol\mu\cdot\hat{\mathbf r})-\boldsymbol\mu}{r^3}.
\]
Se $\theta$ e l'angolo rispetto all'asse magnetico,
\[
B_r=\frac{2\mu\cos\theta}{r^3},\qquad
B_\theta=\frac{\mu\sin\theta}{r^3}.
\]
Alla superficie $r=R$:
\[
B_{\rm eq}=\frac\mu{R^3},\qquad B_{\rm p}=\frac{2\mu}{R^3}=2B_{\rm eq}.
\]
\textbf{Convenzione:} il campo $3.2\times10^{19}\sqrt{P\dot P}\,$G e usualmente quello \emph{equatoriale}, non polare.
\end{frame}

\begin{frame}[t]{Radiazione del dipolo inclinato}
\small
Il momento $\boldsymbol\mu$ forma un angolo $\alpha$ con l'asse di rotazione:
\[
\boldsymbol\mu(t)=\mu
\bigl(\sin\alpha\cos\Omega t,\sin\alpha\sin\Omega t,\cos\alpha\bigr).
\]
Derivando due volte,
\[
|\ddot{\boldsymbol\mu}|^2=\mu^2\Omega^4\sin^2\alpha.
\]
La formula del dipolo magnetico radiativo in vuoto (Gaussiane) fornisce
\[
\mathcal P_{\rm dip}=\frac{2\mu^2\Omega^4}{3c^3}\sin^2\alpha.
\]
Il fattore $\sin^2\alpha$ deriva dalla sola componente perpendicolare all'asse di rotazione.
\end{frame}

\begin{frame}[t]{Derivazione di $B$ da $P$ e $\dot P$}
\small
Assumendo rotazione rigida con $I$ costante e \emph{tutta} la perdita rotazionale in radiazione del dipolo in vuoto:
\[
I\Omega|\dot\Omega|=\frac{2B_{\rm eq}^2R^6\Omega^4}{3c^3}\sin^2\alpha.
\]
Quindi
\[
B_{\rm eq}^2=\frac{3c^3I|\dot\Omega|}{2R^6\Omega^3\sin^2\alpha}.
\]
Usando $\Omega=2\pi/P$, $|\dot\Omega|=2\pi\dot P/P^2$:
\[
B_{\rm eq}=\sqrt{\frac{3c^3I}{8\pi^2R^6}}
\frac{\sqrt{P\dot P}}{|\sin\alpha|}.
\]
Il campo non e una misura diretta: dipende da $I,R,\alpha$ e dal modello del torque.
\end{frame}

\begin{frame}[t]{Campo caratteristico: normalizzazione}
\small
Con $I=10^{45}\,\mathrm{g\,cm^2}$, $R=10^6\,$cm e $\sin\alpha=1$:
\[
B_{\rm char,eq}=3.2\times10^{19}\sqrt{P\dot P}\ \mathrm G
\]
con $P$ in secondi e $\dot P$ in s/s.
\[
B_{\rm p}=2B_{\rm eq}\quad\text{per un dipolo ideale}.
\]
\begin{itemize}
\item Un'inclinazione piu piccola implica $B_{\rm eq}>B_{\rm char,eq}$ nel modello di dipolo \emph{in vuoto}.
\item In una magnetosfera con plasma le coppie e le perdite hanno una diversa dipendenza da $\alpha$.
\item $B_{\rm char}$ vincola soprattutto il campo dipolare a grande scala, non necessariamente i multipoli intensi vicino ai poli.
\end{itemize}
\end{frame}

\begin{frame}[t]{Dipolo nel vuoto vs magnetosfera con plasma}
\small
\begin{columns}[T,onlytextwidth]
\column{.5\textwidth}
\textbf{Dipolo in vuoto}
\[
\mathcal P_{\rm vac}=\frac{2\mu^2\Omega^4}{3c^3}\sin^2\alpha.
\]
\begin{itemize}
\item Predice perdita nulla per $\alpha=0$.
\item Non descrive correnti e vento di particelle.
\end{itemize}
\column{.5\textwidth}
\textbf{Magnetosfera force-free (stima)}
\[
\mathcal P_{\rm ff}\sim\frac{\mu^2\Omega^4}{c^3}(1+\sin^2\alpha).
\]
\begin{itemize}
\item Torque non nullo anche per rotatore allineato.
\item Le correnti magnetosferiche modificano la stima del campo.
\end{itemize}
\end{columns}
\medskip
La dipendenza $\propto\Omega^4$ resta caratteristica, ma il coefficiente non e universale.
\end{frame}

\subsection{6.5.3 e 6.5.4 -- Eta e braking index}

\begin{frame}[t]{Legge di frenamento generale}
\small
Parametrizzare il torque come
\[
\dot\Omega=-K\Omega^n,\qquad K>0,
\]
se $K$ e $n$ costanti nel tempo.
\begin{itemize}
\item $n=3$: dipolo magnetico ideale con $B$, $I$, $\alpha$ costanti.
\item $n=5$: quadrupolo gravitazionale persistente, con ellitticita costante.
\item $n$ osservato in generale misura l'evoluzione \emph{effettiva} del torque.
\end{itemize}
Per $n\neq1$:
\[
\Omega^{1-n}-\Omega_0^{1-n}=(n-1)Kt.
\]
\end{frame}

\begin{frame}[t]{Eta da una legge di torque a potenza}
\small
Dalla legge $\dot\Omega=-K\Omega^n$:
\[
\frac{d\Omega}{\Omega^n}=-K\,dt
\Rightarrow
\frac{\Omega^{1-n}-\Omega_0^{1-n}}{1-n}=-Kt.
\]
Isolando $t$ e sostituendo $K=|\dot\Omega|/\Omega^n$:
\[
t=\frac{\Omega}{(n-1)|\dot\Omega|}
\left[1-\left(\frac{\Omega}{\Omega_0}\right)^{n-1}\right].
\]
Poiche $\Omega/\Omega_0=P_0/P$:
\[
t=\frac{P}{(n-1)\dot P}
\left[1-\left(\frac{P_0}{P}\right)^{n-1}\right].
\]
Vale \emph{solo} se il modello di frenamento resta costante nel tempo.
\end{frame}

\begin{frame}[t]{Eta caratteristica: origine e limiti}
\small
Per $n=3$ e $P_0\ll P$, la formula si riduce a
\[
\tau_c\equiv\frac{P}{2\dot P}.
\]
\begin{itemize}
\item Non coincide necessariamente con l'eta cronologica (cambiamenti di $B,\alpha,I$, coppie esterne, riciclaggio, $P_0$ non trascurabile).
\item Per pulsar binarie $\dot P$ apparente puo includere accelerazioni lungo la linea di vista.
\item Il confronto con l'eta dei resti di supernova verifica ipotesi sul torque e sul periodo iniziale.
\end{itemize}
\[
\tau_c\propto\frac{P}{\dot P},\quad
\text{mentre}\quad B_{\rm char}\propto\sqrt{P\dot P}.
\]
\end{frame}

\begin{frame}[t]{Braking index: definizione osservativa corretta}
\small
Differenziamo $\dot\Omega=-K\Omega^n$ con $K$ costante:
\[
\ddot\Omega=-Kn\Omega^{n-1}\dot\Omega
=n\frac{\dot\Omega^2}{\Omega}.
\]
Segue
\[
n_{\rm br}=\frac{\Omega\ddot\Omega}{\dot\Omega^2}
=\frac{\nu\ddot\nu}{\dot\nu^2}.
\]
In funzione del periodo:
\[
n_{\rm br}=2-\frac{P\ddot P}{\dot P^2}.
\]
\textbf{Correzione:} il segno meno davanti a $\Omega\ddot\Omega/\dot\Omega^2$ riportato nella (6.28) del libro e incompatibile con la sua (6.27).
\end{frame}

\begin{frame}[t]{Braking index diverso da 3: da cosa dipende?}
\small
Se la coppia elettromagnetica effettiva e $\dot\Omega=-K(t)\Omega^3$, allora
\[
\ddot\Omega=-\dot K\Omega^3-3K\Omega^2\dot\Omega.
\]
Dalla definizione osservativa,
\[
n_{\rm br}=3+\frac{\Omega}{\dot\Omega}\frac{\dot K}{K}.
\]
Per il dipolo in vuoto $K\propto B^2R^6\sin^2\alpha/I$:
\[
\frac{\dot K}{K}=2\frac{\dot B}{B}
+6\frac{\dot R}{R}+2\dot\alpha\cot\alpha-\frac{\dot I}{I}.
\]
In pratica possono contribuire anche correnti magnetosferiche, vento e cambiamenti non-secolari dovuti a glitch.
\end{frame}

\begin{frame}[t]{Esempio: il valore del Crab e il modello ideale}
\small
Il testo cita per il Crab un indice effettivo inferiore a $3$, come accade in molte pulsar giovani per cui e stato misurato $\ddot\nu$.
\[
n_{\rm dip}=3,\qquad n_{\rm GW}=5
\quad\text{(modelli ideali a coefficienti costanti).}
\]
\begin{itemize}
\item Un indice $n<3$ suggerisce che il frenamento reale non sia descritto da un dipolo di vuoto con parametri immutabili.
\item Non permette, da solo, di attribuire l'effetto univocamente all'evoluzione del campo magnetico.
\item La misura richiede serie temporali estese: glitch, rumore di timing e cambiamenti di stato possono contaminare $\ddot\nu$.
\end{itemize}
\end{frame}

\subsection{6.5.5 -- Onde gravitazionali continue}

\begin{frame}[t]{Quadrupolo gravitazionale da una stella non assialsimmetrica}
\small
Per una deformazione equatoriale persistente definire l'ellitticita
\[
\epsilon_{\rm ell}\equiv\frac{I_{xx}-I_{yy}}{I_{zz}}.
\]
Per rotazione attorno a $z$, l'emissione dominante e a
\[
f_{\rm GW}=2\nu=\frac2P.
\]
Nell'approssimazione quadrupolare:
\[
\mathcal P_{\rm GW}=\frac{32G}{5c^5}I^2\epsilon_{\rm ell}^2\Omega^6.
\]
\begin{itemize}
\item Un oggetto perfettamente assialsimmetrico con asse di rotazione e simmetria coincidenti non produce questo segnale.
\item La dipendenza $\Omega^6$ e piu ripida di quella dipolare $\Omega^4$.
\end{itemize}
\end{frame}

\begin{frame}[t]{Braking index gravitazionale $n=5$}
\small
Con $I$ costante, uguagliando la potenza persa in rotazione con quella gravitazionale:
\[
I\Omega|\dot\Omega|=\frac{32G}{5c^5}I^2\epsilon_{\rm ell}^2\Omega^6.
\]
Quindi
\[
\dot\Omega=-\underbrace{\frac{32GI\epsilon_{\rm ell}^2}{5c^5}}_{K_{\rm GW}}\Omega^5.
\]
Se $I$ e $\epsilon_{\rm ell}$ sono costanti,
\[
n_{\rm GW}=5.
\]
Un frenamento combinato elettromagnetico e gravitazionale non segue, in generale, una sola potenza $\Omega^n$ con $n$ costante.
\end{frame}

\begin{frame}[t]{Spin-down limit per l'ampiezza delle GW}
\small
L'ampiezza intrinseca di un'onda continua emessa a $2\nu$ da un quadrupolo equatoriale e
\[
h_0=\frac{4GI\epsilon_{\rm ell}\Omega^2}{c^4d}.
\]
Se \emph{tutta} la perdita di energia rotazionale fosse in onde gravitazionali, si ottiene il limite
\[
h_0^{\rm sd}=\frac1d\sqrt{\frac{5GI|\dot\nu|}{2c^3\nu}}.
\]
\begin{itemize}
\item E un massimo energetico stimato, non una previsione dell'ampiezza reale.
\item I limiti interferometrici al di sotto di $h_0^{\rm sd}$ mostrano che le onde gravitazionali non dominano il frenamento per quella sorgente.
\end{itemize}
\end{frame}

\section{Diagramma P--Pdot, evoluzione e death line}

\subsection{6.6 -- Piano periodo--derivata del periodo}

\begin{frame}[t]{Figura 6.1 Schaffner -- Diagramma originale $P$--$\dot P$}
\small
\begin{figure}[!ht]\includegraphics[trim={0cm 0cm 0 0},clip,keepaspectratio,width=0.95\textwidth]{lorimer08/PPdot}\label{fig:lorimer08_PPdot}\end{figure}
\vspace{-1mm}\footnotesize Riproduzione ritagliata della Figura 6.1 del capitolo; dati e classificazioni riferiti alla pubblicazione citata nel libro (Lorimer 2011).
\end{frame}

\begin{frame}[t]{Rette di campo magnetico costante nel piano log-log}
\small
Con la stima dipolare $B_{\rm char}=C_B\sqrt{P\dot P}$:
\[
P\dot P=\frac{B_{\rm char}^2}{C_B^2}
\quad\Longrightarrow\quad
\dot P=\frac{B_{\rm char}^2}{C_B^2P}.
\]
Prendendo i logaritmi:
\[
\log\dot P=-\log P+2\log(B_{\rm char}/C_B).
\]
\begin{itemize}
\item Nel piano $(\log P,\log\dot P)$ le rette $B_{\rm char}=\mathrm{costante}$ hanno pendenza $-1$.
\item In alto a destra dominano generalmente oggetti con campi caratteristici elevati; in basso a sinistra pulsar riciclate.
\end{itemize}
\end{frame}

\begin{frame}[t]{Rette di eta caratteristica costante}
\small
Dalla definizione $\tau_c=P/(2\dot P)$:
\[
\dot P=\frac P{2\tau_c}
\quad\Longrightarrow\quad
\log\dot P=\log P-\log(2\tau_c).
\]
\begin{itemize}
\item Le rette $\tau_c=\mathrm{costante}$ hanno pendenza $+1$ nel piano log-log.
\item Spostarsi verso regioni con $\dot P$ molto bassa, a $P$ fissato, corrisponde a eta caratteristica piu grande.
\item L'eta \emph{caratteristica} non e necessariamente l'eta vera, specialmente per pulsar riciclate.
\end{itemize}
\[
B_{\rm char}\propto(P\dot P)^{1/2},\qquad
\tau_c\propto P/\dot P.
\]
\end{frame}

\begin{frame}[t]{Evoluzione ideale nel diagramma}
\small
Se $n=3$ e $K$ costante:
\[
\dot\Omega=-K\Omega^3
\quad\Longrightarrow\quad P\dot P=4\pi^2K=\mathrm{costante}.
\]
La traiettoria segue quindi una linea di \emph{campo caratteristico costante}, di pendenza $-1$:
\[
\dot P\propto P^{-1}.
\]
\begin{itemize}
\item La pulsar rallenta: $P$ cresce e $\dot P$ decresce lungo la stessa retta.
\item Accrescimento di materia e momento angolare in un sistema binario puo spostare la stella verso piccoli $P$ (riciclaggio).
\item Evoluzione di $B$ e geometria magnetosferica modificano i percorsi.
\end{itemize}
\end{frame}

\begin{frame}[t]{Esempio di diagramma con linea di morte indicativa}
\small
\begin{center}
    \begin{figure}[!ht]\includegraphics[trim={0cm 0cm 0 0},clip,keepaspectratio,width=0.95\textwidth]{astrobites.org/deathline-a_pulsar_alone}\label{fig:astrobites.org_deathline-a_pulsar_alone}\end{figure}
\end{center}
\vspace{-2mm}\footnotesize Grafico didattico creato per gli appunti. La \emph{death line} mostrata e solo una soglia di potenziale illustrativa, non un confine empirico universale.
\end{frame}

\subsection{Approfondimento -- Pulsar death line}

\begin{frame}[t]{Quando una pulsar radio puo spegnersi?}
\small
L'emissione radio coerente e associata a plasma e coppie $e^-e^+$ nelle regioni di accelerazione magnetosferiche.
\begin{itemize}
\item Per oggetti che rallentano, le differenze di potenziale accessibili possono diminuire.
\item Se la cascata di coppie non e piu efficiente, alcuni modelli prevedono la cessazione dell'emissione radio ordinaria.
\item La \emph{pulsar death line} e un confine teorico \emph{dipendente dal modello} nel piano $P$--$\dot P$.
\item Non va confusa con un limite sulla stabilita della stella di neutroni; riguarda il meccanismo radiativo.
\end{itemize}
\end{frame}

\begin{frame}[t]{Potenziale dei poli e sua dipendenza da $P,B$}
\small
La rotazione e il campo magnetico producono una scala di potenziale ai poli:
\[
\Delta\Phi_{\rm pc}\sim\frac{B_{\rm p}R^3\Omega^2}{2c^2}
=\frac{2\pi^2B_{\rm p}R^3}{c^2P^2}.
\]
Quindi, a $R$ fissato,
\[
\Delta\Phi_{\rm pc}\propto\frac{B}{P^2}.
\]
Un criterio modellistico schematizzato come $\Delta\Phi_{\rm pc}\gtrsim\Phi_{\rm crit}$ equivale a
\[
\frac B{P^2}\gtrsim\mathrm{costante}.
\]
La produzione effettiva di coppie dipende anche da curvatura delle linee, fotoni seed, moltipolarita e campi elettrici paralleli.
\end{frame}

\begin{frame}[t]{Death line: derivazione della pendenza $+3$}
\small
Nel modello elementare del dipolo:
\[
B=C_B(P\dot P)^{1/2}.
\]
Alla soglia di potenziale, $B/P^2=C_\Phi$:
\[
\frac{C_B(P\dot P)^{1/2}}{P^2}=C_\Phi
\quad\Rightarrow\quad
\dot P=\left(\frac{C_\Phi}{C_B}\right)^2P^3.
\]
Pertanto
\[
\log\dot P=3\log P+\mathrm{costante}.
\]
\begin{itemize}
\item Questa pendenza $+3$ e una \emph{semplificazione}, non una derivazione microscopica universale della death line.
\item Soglie basate su fotoni di curvatura, Compton inverso e campi multipolari producono confini differenti (o vere \emph{death valleys}).
\end{itemize}
\end{frame}

\begin{frame}[t]{Death line vs eta e campo: non confonderle}
\small
\begin{tabular}{@{}lll@{}}
\toprule
Condizione & Legge nel piano $P$--$\dot P$ & Pendenza \\
\midrule
$B_{\rm char}$ costante & $\dot P\propto P^{-1}$ & $-1$ \\
$\tau_c$ costante & $\dot P\propto P^{+1}$ & $+1$ \\
Potenziale polare di soglia & $\dot P\propto P^{+3}$ & $+3$ \\
\bottomrule
\end{tabular}
\medskip
\begin{itemize}
\item Le prime due derivano da definizioni osservativo-modellistiche standard.
\item La terza deriva da un criterio di morte \emph{ipotetico e semplificato}.
\item Una pulsar sotto una linea disegnata come death line non costituisce automaticamente un paradosso: e il modello della cascata da verificare.
\end{itemize}
\end{frame}

\section{Pulsar glitches e superfluidita}

\subsection{6.7 -- Salti nel periodo}

\begin{frame}[t]{Che cos'e un glitch?}
\small
Il timing mostra talvolta un aumento improvviso della frequenza rotazionale:
\[
\Delta\nu>0,\qquad
\frac{\Delta\nu}{\nu}\sim10^{-11}\text{--}10^{-5}
\quad\text{(intervallo qualitativo del testo).}
\]
Con $P=1/\nu$, al primo ordine
\[
\frac{\Delta P}{P}\simeq-\frac{\Delta\nu}{\nu}<0.
\]
\begin{itemize}
\item Dopo il glitch possono seguire un recupero transitorio e un cambiamento della velocita di frenamento.
\item La dinamica coinvolge il momento angolare trasferito fra componenti interne con rotazioni diverse.
\end{itemize}
\end{frame}

\begin{frame}[t]{Due-componenti: crosta e superfluido}
\small
Modello schematico: componente osservata $c$ (crosta + parte accoppiata) e componente superfluida $s$.
\[
J=I_c\Omega_c+I_s\Omega_s.
\]
Durante un glitch rapido, trascurando il torque esterno:
\[
\Delta J\simeq I_c\Delta\Omega_c+I_s\Delta\Omega_s=0.
\]
Percio
\[
\Delta\Omega_c=-\frac{I_s}{I_c}\Delta\Omega_s.
\]
Se il superfluido rallenta ($\Delta\Omega_s<0$), la crosta accelera ($\Delta\Omega_c>0$).
\medskip
Questa equazione mostra come i glitch possano vincolare la \emph{frazione di momento d'inerzia} del serbatoio superfluido.
\end{frame}


\begin{frame}[t]{Haensel \S9.7: il serbatoio di momento angolare}
\small
Il modello distingue una componente accoppiata alla crosta, $I_c$, e un serbatoio superfluido, $I_{\rm res}$.
La rotazione della crosta rallenta sotto il torque esterno; il superfluido debolmente accoppiato conserva temporaneamente una velocita' maggiore.
\[
\delta\Omega=\Omega_{\rm res}-\Omega_c>0.
\]
Durante il glitch, su tempi brevi, il trasferimento interno soddisfa
\[
I_c\Delta\Omega_c+I_{\rm res}\Delta\Omega_{\rm res}\simeq0.
\]
L'unpinning dei vortici consente di trasferire momento angolare al componente osservabile.
\par\medskip
\footnotesize Haensel et al. (2006), \S9.7, pp.~508--509.
\end{frame}

\begin{frame}[t]{Haensel \S9.7: derivazione del vincolo dai glitches}
\small
Definiamo l'attivita' osservata dei glitches in un intervallo $T$:
\[
\mathcal A=\frac1T\sum_i\frac{\Delta\Omega_i}{\Omega}.
\]
Il momento angolare cumulativo trasferito e'
\[
\Delta J_{\rm obs}=I_c\sum_i\Delta\Omega_i
=I_c\Omega\mathcal A T.
\]
Il serbatoio non puo' accumulare piu' momento angolare di quello generato dal rallentamento differenziale:
\[
\Delta J_{\rm obs}\le I_{\rm res}|\dot\Omega|T
\quad\Longrightarrow\quad
\frac{I_{\rm res}}{I_c}\ge\mathcal A\frac{\Omega}{|\dot\Omega|}.
\]
\footnotesize Haensel et al. (2006), eqs.~(9.43)--(9.47); ipotesi di accoppiamento del modello.
\end{frame}

\begin{frame}[t]{Haensel \S9.7: Vela e il contributo della crosta}
\small
Per Vela il testo riporta, con dati disponibili fino al 2005,
\[
\frac{I_{\rm res}}{I_c}\gtrsim 1.5\%.
\]
Se il serbatoio fosse interamente nella crosta interna, il vincolo diventerebbe una condizione su $I_{\rm crust}/I$.
\begin{itemize}
\item Non e' una misura diretta del momento d'inerzia crostale.
\item Dipende dall'ipotesi che il serbatoio sia crostale e dalla dinamica di accoppiamento.
\item L'entrainment dei neutroni e il possibile contributo del core modificano l'interpretazione quantitativa.
\end{itemize}
\footnotesize Haensel et al. (2006), \S9.7, pp.~510--511. Il valore e' storico, non una stima aggiornata.
\end{frame}

\begin{frame}[t]{Vortici quantizzati nella materia neutronica}
\small
Per coppie di neutroni di massa circa $2m_n$, la circolazione della velocita superfluida e quantizzata:
\[
\oint\mathbf v_s\cdot d\boldsymbol\ell=
\kappa\,N_v,\qquad
\kappa=\frac h{2m_n}.
\]
Per rotazione uniforme, relazione di Feynman:
\[
n_v\kappa=2\Omega_s
\quad\Rightarrow\quad
n_v=\frac{2\Omega_s}{\kappa}.
\]
\begin{itemize}
\item Vortici \emph{pinnati} al reticolo nucleare possono trattenere momento angolare.
\item Unpinning e moto collettivo trasferiscono momento angolare alla componente osservata.
\item \emph{Entrainment} e accoppiamento crosta-core modificano la frazione $I_s$ effettivamente disponibile.
\end{itemize}
\end{frame}

\begin{frame}[t]{Glitches: collegamento con la teoria BCS}
\small
Il pairing neutronico discusso nel capitolo 3 e una spiegazione microscopica della comparsa del superfluido.
\begin{itemize}
\item Nella crosta interna prevale il pairing $^1S_0$ dei neutroni.
\item Nelle regioni piu dense si considera anche pairing anisotropo, per esempio $^3P_2$.
\item L'ordine di fase macroscopico permette vortici quantizzati, assenti in una descrizione di fluido normale ideale.
\end{itemize}
\[
\text{pairing BCS}\ \Longrightarrow\ \text{superfluidita}
\ \Longrightarrow\ \text{vortici}\ \Longrightarrow\ \text{possibili glitch}.
\]
Non ogni glitch identifica un unico meccanismo microscopico: sono importanti anche elasticita, campi magnetici e trasporto.
\end{frame}

\section{Magnetosfera del rotatore allineato}

\subsection{6.8.1 e 6.8.2 -- Campo elettrico e cilindro di luce}

\begin{frame}[t]{Conduttore in rotazione: campo indotto}
\small
Nel limite di conduttivita ideale in un riferimento che coruota:
\[
\mathbf E+\frac{\mathbf v}{c}\times\mathbf B=0,
\qquad \mathbf v=\boldsymbol\Omega\times\mathbf r
\quad\text{(Gaussiane).}
\]
Per $\boldsymbol\Omega\parallel\mathbf B=B_0\hat z$, $\mathbf v=\Omega r\sin\theta\,\hat\varphi$:
\[
\mathbf E=-\frac{\mathbf v}{c}\times\mathbf B.
\]
\begin{itemize}
\item Il campo elettrico interno non e nullo nel riferimento inerziale: compensa la forza magnetica sulle cariche corotanti.
\item La relazione si applica come condizione di idealita; la magnetosfera reale puo avere regioni di accelerazione con $E_\parallel\neq0$.
\end{itemize}
\end{frame}

\begin{frame}[t]{Potenziale elettrostatico del rotatore allineato}
\small
In un campo uniforme $\mathbf B=B_0\hat z$:
\[
\mathbf v\times\mathbf B
=\Omega B_0 r\sin\theta
\bigl(\sin\theta\,\hat r+\cos\theta\,\hat\theta\bigr)
\]
con convenzione di verso scelta coerentemente con $\Omega,B_0>0$.
\[
\mathbf E=-\frac{\Omega B_0}{c}r\sin\theta
\bigl(\sin\theta\,\hat r+\cos\theta\,\hat\theta\bigr).
\]
Un potenziale che soddisfa $\mathbf E=-\boldsymbol\nabla\Phi$ e
\[
\Phi(r,\theta)=\frac{\Omega B_0}{2c}r^2\sin^2\theta+\text{cost}.
\]
Il segno dipende dall'orientazione relativa di $\mathbf B$ e $\boldsymbol\Omega$; le differenze di potenziale fisiche contano.
\end{frame}

\begin{frame}[t]{Cilindro di luce: derivazione}
\small
La velocita di corotazione a distanza cilindrica $\varpi$ dall'asse e
\[
v_\varphi=\Omega\varpi.
\]
Il valore formale $v_\varphi=c$ determina il \emph{light cylinder}:
\[
R_L=\frac c\Omega=\frac{cP}{2\pi}.
\]
Per $P=1\,$s:
\[
R_L\simeq4.77\times10^4\,\mathrm{km}.
\]
Per $P=1\,$ms:
\[
R_L\simeq47.7\,\mathrm{km}.
\]
\textbf{Nota:} il light cylinder e una superficie cilindrica rispetto all'asse di rotazione, non un raggio sferico di confinamento della materia.
\end{frame}

\begin{frame}[t]{Linee di campo di un dipolo}
\small
Per un campo magnetico meridionale $\mathbf B=B_r\hat r+B_\theta\hat\theta$, la tangente alla linea di campo soddisfa
\[
\frac{dr}{r\,d\theta}=\frac{B_r}{B_\theta}
=2\cot\theta.
\]
Integrando:
\[
\frac{dr}{r}=2\cot\theta\,d\theta
\quad\Rightarrow\quad
\ln r=2\ln\sin\theta+\ln C.
\]
Percio
\[
r=C\sin^2\theta.
\]
Nel modello dipolare, le linee polari che si estendono oltre il cilindro di luce formano una regione \emph{aperta}.
\end{frame}

\begin{frame}[t]{Calotta polare: angolo di apertura}
\small
La linea limite tocca il light cylinder presso l'equatore: $C\simeq R_L$.
Alla superficie $r=R$:
\[
R=R_L\sin^2\theta_{\rm pc}
\quad\Rightarrow\quad
\sin\theta_{\rm pc}\simeq\sqrt{\frac R{R_L}}.
\]
Per $R\ll R_L$, $\theta_{\rm pc}\ll1$, dunque
\[
r_{\rm pc}=R\sin\theta_{\rm pc}
\simeq R\sqrt{\frac R{R_L}}
=R\sqrt{\frac{\Omega R}c}.
\]
Per $R=10\,$km e $P=1\,$s:
\[
r_{\rm pc}\simeq145\,\mathrm m.
\]
Questa e una stima geometrica per un dipolo allineato.
\end{frame}

\subsection{6.8.3 e 6.8.4 -- Accelerazione e Poynting flux}

\begin{frame}[t]{Potenziale disponibile ai poli}
\small
Dalla soluzione elettrostatica e dalla dimensione della calotta:
\[
\Delta\Phi_{\rm pc}\simeq\frac{\Omega B_0 R^2}{2c}
\sin^2\theta_{\rm pc}.
\]
Poiche $\sin^2\theta_{\rm pc}\simeq R/R_L=\Omega R/c$:
\[
\Delta\Phi_{\rm pc}\simeq
\frac{B_0\Omega^2R^3}{2c^2}.
\]
La scala di energia per una carica $e$ e
\[
E_{\rm max,ideal}\sim e\Delta\Phi_{\rm pc}.
\]
In pratica accelerazione e scariche di coppie limitano i potenziali \emph{non schermati} disponibili.
\end{frame}

\begin{frame}[t]{Generazione di coppie e plasma magnetosferico}
\small
\begin{itemize}
\item Campi elettrici paralleli alle linee di campo possono accelerare elettroni o positroni relativistici.
\item La radiazione di curvatura e i processi di Compton inverso producono fotoni energetici.
\item In campo magnetico forte, fotoni possono creare coppie $e^+e^-$ e avviare una cascata.
\item La moltiplicazione di cariche tende a schermare $E_\parallel$ e modifica la struttura globale del torque.
\end{itemize}
\[
\text{accelerazione}\to\gamma\to e^+e^-\to\text{plasma magnetosferico}.
\]
Il legame con la \emph{death line} e il venir meno, in alcuni regimi, di una cascata sufficiente.
\end{frame}

\begin{frame}[t]{Energia elettromagnetica in uscita}
\small
Il flusso di Poynting in unita gaussiane e
\[
\mathbf S=\frac c{4\pi}\mathbf E\times\mathbf B.
\]
Vicino al cilindro di luce, per ordine di grandezza $E\sim B$ e
\[
\mathcal P\sim4\pi R_L^2\frac{cB_L^2}{4\pi}
\sim cR_L^2B_L^2.
\]
Usando il dipolo $B_L\sim B_0(R/R_L)^3$:
\[
\mathcal P\sim cB_0^2R^6/R_L^4
\sim\frac{B_0^2R^6\Omega^4}{c^3}.
\]
Si recupera la dipendenza $\Omega^4$, anche senza un modello di emissione dipolare nel vuoto.
\end{frame}

\section{Propagazione radio e dispersion measure}

\subsection{6.9 -- Plasma interstellare}

\begin{frame}[t]{Dispersion measure: definizione fisica}
\small
Un impulso radio attraversa plasma interstellare con densita elettronica $n_e(s)$.
La misura della colonna di elettroni e
\[
\mathrm{DM}=\int_0^D n_e(s)\,ds.
\]
Unita astronomiche standard: $\mathrm{pc\,cm^{-3}}$.
\begin{itemize}
\item DM non e una distanza misurata direttamente.
\item Per ottenere $D$, occorre un modello della distribuzione elettronica lungo la linea di vista o altri vincoli (parallasse, compagna, ecc.).
\item La dispersione differenziale fra frequenze permette di misurare DM con precisione dai tempi di arrivo.
\end{itemize}
\end{frame}

\begin{frame}[t]{Frequenza di plasma: oscillatore di cariche}
\small
Nel sistema di unita gaussiane, per uno spostamento $x$ degli elettroni rispetto a un fondo ionico, il campo di richiamo da
\[
m_e\ddot x=-4\pi e^2n_e x.
\]
Quindi
\[
\ddot x+\omega_p^2 x=0,\qquad
\omega_p^2=\frac{4\pi e^2n_e}{m_e}.
\]
In SI la stessa frequenza e $\omega_p^2=n_e e^2/(\varepsilon_0m_e)$.
\medskip
\textbf{Attenzione alle unita:} e scorretto combinare direttamente i fattori $4\pi$ gaussiani con $\varepsilon_0$ del SI.
\end{frame}

\begin{frame}[t]{Relazione di dispersione e velocita di gruppo}
\small
Per onde elettromagnetiche trasverse nel plasma freddo non magnetizzato:
\[
\omega^2=\omega_p^2+c^2k^2.
\]
Quindi
\[
v_g=\frac{d\omega}{dk}=\frac{c^2k}{\omega}
=c\sqrt{1-\frac{\omega_p^2}{\omega^2}}.
\]
Se $\omega\gg\omega_p$, sviluppando la radice:
\[
v_g\simeq c\left(1-\frac{\omega_p^2}{2\omega^2}\right).
\]
Le frequenze radio piu basse subiscono quindi un ritardo piu grande.
\end{frame}

\begin{frame}[t]{Derivazione del ritardo radio $\Delta t\propto\mathrm{DM}\,\nu^{-2}$}
\small
Dal tempo di propagazione:
\[
t(\omega)=\int_0^D\frac{ds}{v_g(s,\omega)}
\simeq\frac Dc+\frac1{2c\omega^2}\int_0^D\omega_p^2(s)ds.
\]
Inserendo $\omega_p^2=4\pi e^2n_e/m_e$:
\[
t(\nu)-D/c\propto\frac{\mathrm{DM}}{\nu^2}.
\]
Formula pratica del ritardo \emph{fra due frequenze}, con $\nu$ in GHz:
\[
\Delta t\,[\mathrm{ms}]
\simeq4.149\,\mathrm{DM}\,[\mathrm{pc\,cm^{-3}}]
\left(\nu_1^{-2}-\nu_2^{-2}\right).
\]
\end{frame}

\begin{frame}[t]{Esempio numerico di dispersione}
\small
Per $\mathrm{DM}=100\,\mathrm{pc\,cm^{-3}}$ e osservazioni a $\nu_1=1.0\,$GHz, $\nu_2=1.5\,$GHz:
\[
\begin{aligned}
\Delta t&\simeq4.149\times100\left(1-\frac1{1.5^2}\right)\,\mathrm{ms}\\
&\simeq230.5\,\mathrm{ms}.
\end{aligned}
\]
\begin{itemize}
\item Questa differenza e molto superiore alla durata tipica di molti impulsi.
\item La \emph{dedispersion} corregge il ritardo in frequenza prima della ricerca e della misura precisa dei TOA.
\item DM puo variare se cambia la colonna elettronica lungo il cammino.
\end{itemize}
\end{frame}

\section{Pulsar binarie: masse e prove di relativita}

% Integrazione didattica da Haensel, Potekhin, Yakovlev (2006), §§9.1--9.2.
% Le masse riportate come esempi storici si riferiscono al 2005--2006.
\subsection{Haensel 9.1 -- Metodi di misura delle masse}

\begin{frame}[t]{Masse osservate e vincoli sull'equazione di stato}
\small
\begin{itemize}
\item Le masse inferite dalle orbite sono \emph{masse gravitazionali}, non masse barioniche.
\item Una misura precisa di $M$ vincola la curva $M(R)$ prevista dall'EOS; una massa elevata impone $M_{\max}^{\rm EOS}\geq M_{\rm osservata}$ (per configurazioni confrontabili quanto a rotazione).
\item Sistemi osservativi: binarie X con compagna ottica, doppie stelle di neutroni, pulsar--nana bianca, pulsar--stella non degenere.
\item Accuratezza statistica e incertezze sistematiche sono diverse: il timing delle pulsar compatte e spesso piu robusto delle velocita radiali di stelle estese.
\end{itemize}
\medskip\footnotesize Fonte: Haensel et al., \S9.1, pp.~456--489. Le statistiche della figura 9.1 sono riferite al 2005.
\end{frame}

\begin{frame}[t]{Geometria orbitale e velocita radiale}
\small
Per $M=M_p+M_c$, $a=a_p+a_c$, $n_b=2\pi/P_b$:
\[
M_pa_p=M_ca_c,\qquad n_b^2a^3=GM.
\]
Per una componente osservata spettroscopicamente,
\[
v_r(t)=K\,[\cos(\omega+\phi(t))+e\cos\omega],\qquad
K=\frac{n_ba_p\sin i}{\sqrt{1-e^2}}.
\]
\begin{itemize}
\item Lo spostamento Doppler delle righe ottiche misura $K$; il timing X/radio misura $x=a_p\sin i/c$.
\item $P_b,e,\omega,T_0,x$ descrivono l'orbita kepleriana della componente osservata.
\item La velocita radiale non determina da sola $i$ e le due masse.
\end{itemize}
\end{frame}

\begin{frame}[t]{Funzione di massa: derivazione e degenerazione}
\small
Dal baricentro $a_p=(M_c/M)a$ e da Keplero $n_b^2a^3=GM$:
\[
f_p\equiv\frac{n_b^2(a_p\sin i)^3}{G}
=\frac{(M_c\sin i)^3}{M^2}
=\frac{4\pi^2c^3x^3}{GP_b^2}.
\]
\begin{itemize}
\item $f_p$ e direttamente osservabile, ma dipende da $M_p,M_c,i$.
\item Siccome $\sin i\leq1$ e $M\geq M_c$, segue $M_c\geq f_p$; con $M_p$ fissata si ottiene un limite inferiore piu stringente.
\item Se si misurano entrambe le velocita: $M_p/M_c=K_c/K_p=a_c/a_p$; resta necessaria una misura di inclinazione o un ulteriore vincolo.
\end{itemize}
\end{frame}

\begin{frame}[t]{Binarie X: come si chiude il sistema}
\small
\begin{itemize}
\item Timing delle pulsazioni X: orbita della stella compatta ($K_p$ o $x$).
\item Spettroscopia ottica: orbita della compagna ($K_c$), da cui $q=M_p/M_c=K_c/K_p$.
\item Eclissi, curva di luce e geometria del lobo di Roche: vincoli su $i$ e sul raggio della compagna.
\item Allargamento rotazionale delle righe: vincolo addizionale se la rotazione della compagna e sincronizzata.
\item \emph{Limite sistematico}: venti, accrescimento, oscillazioni e deformazioni mareali possono alterare la curva di velocita radiale.
\end{itemize}
\medskip\footnotesize Esempio storico: Vela X-1, Haensel et al., \S9.1.1b--c; stime di massa dipendenti dalle ipotesi geometriche.
\end{frame}

\begin{frame}[t]{Pulsar binarie: timing e parametri post-kepleriani}
\small
Il modello di tempi di arrivo separa diversi ritardi:
\[
\Delta t=\Delta_{\rm R}+\Delta_{\rm E}+\Delta_{\rm S}+\cdots
\]
\begin{itemize}
\item Roemer $\Delta_{\rm R}$: propagazione geometrica attraverso l'orbita.
\item Einstein $\Delta_{\rm E}$: Doppler trasverso e redshift gravitazionale \emph{variabili}.
\item Shapiro $\Delta_{\rm S}$: propagazione nel campo gravitazionale della compagna.
\item Derivate secolari $\dot\omega$, $\dot P_b$: precessione relativistica e perdita di energia orbitale.
\end{itemize}
\medskip
In relativita generale due parametri post-kepleriani indipendenti possono determinare $M_p,M_c$; gli altri testano la consistenza della teoria.
\end{frame}

\begin{frame}[t]{Periastron: misura della massa totale}
\small
Con $n_b=2\pi/P_b$ e $M=M_p+M_c$:
\[
\dot\omega=\frac{3n_b^{5/3}(GM)^{2/3}}{c^2(1-e^2)},
\qquad
M=\frac1G\left[\frac{\dot\omega c^2(1-e^2)}{3n_b^{5/3}}\right]^{3/2}.
\]
\begin{itemize}
\item Il periastron deve essere osservativamente definibile: quando $e\to0$, la direzione del periastron diventa indeterminata.
\item A $P_b,M$ fissati, il fattore $1/(1-e^2)$ amplifica $\dot\omega$, ma non garantisce da solo una misura precisa.
\item Al livello di precisione elevato si devono considerare eventuali contributi non relativistici o di ordine superiore.
\end{itemize}
\end{frame}

\begin{frame}[t]{Redshift orbitale: dal tempo proprio al timing}
\small
Per la pulsar in campo debole, a ordine $c^{-2}$, schematicamente:
\[
\frac{d\tau_p}{dt}=1-\frac{v_p^2}{2c^2}-\frac{GM_c}{rc^2}
+\text{contributi costanti}+O(c^{-4}).
\]
\begin{itemize}
\item Il termine $v_p^2/(2c^2)$ e la dilatazione temporale cinematica.
\item Il termine $GM_c/(rc^2)$ e il redshift gravitazionale dovuto alla compagna.
\item In un'orbita eccentrica $r(t)$ e $v_p(t)$ variano; integrando la parte periodica si ottiene il \emph{ritardo Einstein}.
\item Un redshift costante e assorbito nella frequenza di rotazione intrinseca stimata dal timing.
\end{itemize}
\end{frame}

\begin{frame}[t]{Parametro Einstein: dipendenza dalle masse}
\small
Con $u$ anomalia eccentrica e $u-e\sin u=n_b(t-T_0)$:
\[
\Delta_{\rm E}(t)=\gamma_{\rm E}\sin u(t),
\]
\[
\gamma_{\rm E}=e\left(\frac{P_b}{2\pi}\right)^{1/3}
\frac{G^{2/3}}{c^2}
\frac{M_c(M_p+2M_c)}{(M_p+M_c)^{4/3}}.
\]
\begin{itemize}
\item $\gamma_{\rm E}$ ha dimensione di tempo; e misurata dalle variazioni periodiche dei tempi di arrivo.
\item $\gamma_{\rm E}\propto e$: in un'orbita circolare questa modulazione scompare.
\item Non e una misura isolata del solo redshift gravitazionale: comprende anche la dilatazione temporale cinematica.
\end{itemize}
\medskip\footnotesize Haensel et al., eq.~(9.22), p.~472.
\end{frame}

\begin{frame}[t]{Due osservabili, due masse: $\dot\omega$ e $\gamma_{\rm E}$}
\small
\begin{enumerate}
\item Dal periodo $P_b$, dall'eccentricita $e$ e da $\dot\omega$ si ricava $M=M_p+M_c$.
\item Sostituendo $M_p=M-M_c$ nel ritardo Einstein:
\[
\gamma_{\rm E}=A\,M_c(2M-M_c),\qquad
A=\frac{eG^{2/3}}{c^2M^{4/3}}\left(\frac{P_b}{2\pi}\right)^{1/3}.
\]
\item Si risolve per $M_c$ scegliendo la soluzione fisica $0<M_c<M$; poi $M_p=M-M_c$.
\end{enumerate}
\medskip
Un terzo osservabile indipendente, per esempio $\dot P_b$, non serve piu a fissare le due masse: diventa un test della relativita generale (dopo le correzioni cinematiche e ambientali).
\end{frame}

\begin{frame}[t]{Ritardo di Shapiro: massa compagna e inclinazione}
\small
Nel timing relativistico si introducono il \emph{range} e la \emph{shape}:
\[
r_{\rm S}=\frac{GM_c}{c^3},\qquad s_{\rm S}=\sin i.
\]
\begin{itemize}
\item Il ritardo e maggiore quando il segnale passa vicino alla compagna, specialmente in sistemi visti quasi di taglio ($i\simeq90^\circ$).
\item $r_{\rm S}$ fornisce $M_c$; $s_{\rm S}$ fornisce $i$.
\item La funzione di massa completa la determinazione di $M_p$.
\item A differenza di $\gamma_{\rm E}$, Shapiro puo essere misurato anche per orbite quasi circolari.
\end{itemize}
\medskip\footnotesize Haensel et al., eq.~(9.23), p.~472. $r_{\rm S}$ non e il raggio stellare.
\end{frame}

\begin{frame}[t]{Perche misurare (o limitare) l'eccentricita?}
\small
\begin{itemize}
\item $e$ determina la forma della curva di velocita radiale e i ritardi di timing; un limite $e<e_{\max}$ e un'informazione osservativa, non un valore misurato nullo.
\item Per $e\to0$ il periastron perde significato geometrico e $\gamma_{\rm E}\to0$.
\item Nelle binarie riciclate con nane bianche, orbite quasi circolari sono compatibili con la circolarizzazione durante il trasferimento di massa.
\item Nelle doppie stelle di neutroni, eccentricita significative possono conservare traccia della storia evolutiva e della seconda supernova.
\item Per orbite quasi circolari sono utili parametrizzazioni di timing con $e\sin\omega$ e $e\cos\omega$ anziche $e,\omega$ separatamente.
\end{itemize}
\end{frame}

\begin{frame}[t]{Onde gravitazionali: evoluzione di $a$ ed $e$}
\small
Nel limite quadrupolare (Peters--Mathews), con $M=M_p+M_c$:
\[
\dot a=-\frac{64G^3M_pM_cM}{5c^5a^3}
\frac{1+\frac{73}{24}e^2+\frac{37}{96}e^4}{(1-e^2)^{7/2}},
\]
\[
\dot e=-\frac{304G^3M_pM_cM}{15c^5a^4}
\frac{e(1+\frac{121}{304}e^2)}{(1-e^2)^{5/2}}.
\]
\begin{itemize}
\item Per $0<e<1$: $\dot a<0$ e $\dot e<0$; l'orbita si restringe e si circolarizza.
\item Per $e=0$, $\dot e=0$; la soluzione circolare va trattata direttamente.
\end{itemize}
\medskip\footnotesize Haensel et al., eqs.~(9.11)--(9.20), pp.~468--470.
\end{frame}

\begin{frame}[t]{Figura 9.4 -- Evoluzione della binaria Hulse--Taylor}
    \begin{figure}[!ht]\includegraphics[trim={0cm 0cm 0 0},clip,keepaspectratio,width=0.95\textwidth]{haensel07/hulse-taylor_orbit}\label{fig:haensel07_hulse-taylor_orbit}\end{figure}
\par\vspace{1mm}\scriptsize Haensel et al. (2006), fig.~9.4, p.~470: ricostruzione storica dell'evoluzione di $a$, $e$, $P_b$ e $\dot\omega$.
\end{frame}

\begin{frame}[t]{Diagramma massa--massa: misure e test della RG}
\small
\begin{itemize}
\item Ogni parametro post-kepleriano misurato definisce una curva (con fascia d'errore) nel piano $(M_p,M_c)$, \emph{assumendo una teoria gravitazionale}.
\item $\dot\omega$ fissa principalmente $M_p+M_c$; $\gamma_{\rm E}$ seleziona una seconda curva.
\item L'intersezione determina le masse; $\dot P_b$ e altri parametri verificano la coerenza.
\item Correzioni al $\dot P_b$ osservato: accelerazioni lungo la linea di vista e moto proprio (effetto Shklovskii), quando rilevanti.
\end{itemize}
\begin{figure}[!ht]\includegraphics[trim={0cm 0cm 0 0},clip,keepaspectratio,width=0.95\textwidth]{haensel07/Mass_HulseTaylor_pulsar}\label{fig:haensel07_Mass_HulseTaylor_pulsar}\end{figure}
\par\scriptsize Haensel et al., fig.~9.5: evoluzione storica delle stime delle due masse di PSR B1913+16 ($2\sigma$).
\end{frame}

\begin{frame}[t]{Doppia pulsar e sistemi pulsar--nana bianca}
\small
\begin{itemize}
\item In J0737--3039A/B il timing di entrambe le pulsar permette di misurare $x_A$ e $x_B$, quindi il rapporto delle masse $M_A/M_B=x_B/x_A$ al primo ordine post-newtoniano rilevante.
\item Ulteriori parametri post-kepleriani rendono il sistema un test sovradeterminato della relativita generale.
\item Nelle pulsar--nana bianca si possono combinare timing, spettroscopia ottica della nana bianca e modelli atmosferici (gravita superficiale, temperatura, raggio).
\item Per sistemi riciclati quasi circolari la relazione evolutiva $P_b$--$M_{\rm WD}$ e un vincolo \emph{modellistico}, non una misura dinamica diretta.
\end{itemize}
\medskip\footnotesize Haensel et al., \S\S9.1.2d e 9.1.3, pp.~479--486.
\end{frame}

\begin{frame}[t]{Masse elevate: cosa escludono dell'EOS?}
\small
\begin{itemize}
\item Per una stella non rotante stabile, una EOS ammissibile deve consentire una massa massima almeno pari alla massa gravitazionale osservata.
\item Una misura con errore statistico piccolo non e necessariamente affidabile se dipende da ipotesi su inclinazione, lobo di Roche o atmosfera.
\item Haensel et al. (2006) presentano un limite storico $M_{\max}^{\rm EOS}>1.6\,M_\odot$ (eq.~9.26): \emph{non va interpretato come limite osservativo attuale}.
\item La presenza di stelle massicce discrimina EOS troppo soffici, ma la massa da sola non ricostruisce univocamente $P(\varepsilon)$.
\end{itemize}
\end{frame}

\subsection{Haensel 9.2 -- Redshift gravitazionale superficiale}

\begin{frame}[t]{Redshift superficiale: dalla metrica alla compattezza}
\small
Per una stella sferica statica, all'esterno vale Schwarzschild:
\[
\frac{\nu_\infty}{\nu_{\rm em}}=\sqrt{1-\frac{2GM}{Rc^2}},\qquad
1+z_s=\left(1-\frac{2GM}{Rc^2}\right)^{-1/2}.
\]
Invertendo:
\[
\boxed{\frac{GM}{Rc^2}=\frac12\left[1-(1+z_s)^{-2}\right]}.
\]
\begin{itemize}
\item Una riga spettrale identificata con energia propria $E_0$ e osservata a $E_\infty=E_0/(1+z_s)$.
\item La misura di $z_s$ determina $M/R$, non $M$ e $R$ separatamente.
\end{itemize}
\end{frame}

\begin{frame}[t]{Redshift assoluto: perche e difficile?}
\small
\begin{itemize}
\item Bisogna identificare con sicurezza la transizione e conoscere la sua energia propria.
\item Doppler orbitale e sistemico, rotazione, campi magnetici, atmosfera e luogo di formazione della riga alterano il segnale.
\item Se la riga si forma a $r_{\rm em}>R$, il redshift gravitazionale misurato e inferiore a quello superficiale: $z(r_{\rm em})<z_s$.
\item Nel timing delle pulsar un redshift \emph{costante} e degenere con la frequenza intrinseca di rotazione; una modulazione orbitale periodica e invece separabile.
\end{itemize}
\medskip
\textbf{Non confondere}: $z_s$ e un rapporto di frequenze spettrali; $\gamma_{\rm E}$ e un'ampiezza temporale che combina gravita e cinematica.
\end{frame}

\begin{frame}[t]{Caso storico: righe nei burst e EXO 0748--676}
\small
\begin{itemize}
\item Haensel et al. discutono le prime interpretazioni di caratteristiche gamma come annichilazione $e^+e^-\to2\gamma$ redshiftata; molte di queste identificazioni non furono confermate.
\item Per SGR 0526--66 si discusse una riga a circa $430\,\mathrm{keV}$, interpretata condizionatamente come $511\,\mathrm{keV}$ redshiftati: $z\sim0.23$.
\item Per EXO 0748--676 venne proposto $z_s\simeq0.35$ da righe di assorbimento X attribuite a Fe e O.
\item \emph{Avvertenza}: si tratta di interpretazioni storiche, non di misure definitive; identificazione e riproducibilita delle righe sono cruciali.
\end{itemize}
\medskip\footnotesize Haensel et al., \S9.2, pp.~490--491.
\end{frame}

\begin{frame}[t]{Redshift, massa e raggio: un vincolo complementare}
\small
Definendo $r_g=2GM/c^2$, dalla misura di $z_s$:
\[
R=\frac{r_g(1+z_s)^2}{z_s(2+z_s)}.
\]
Per esempio, se $M=1.4M_\odot$ e $R=12\,\mathrm{km}$,
\[
\frac{2GM}{Rc^2}\simeq0.345,\qquad z_s\simeq0.236.
\]
\begin{itemize}
\item Un vincolo indipendente sulla massa converte il redshift in un vincolo sul raggio.
\item Una coppia $(M,R)$ confrontata con le curve TOV puo discriminare le EOS.
\item La formula assume simmetria sferica e assenza di rotazione significativa; per stelle rapidamente rotanti occorre un modello piu completo.
\end{itemize}
\medskip\footnotesize Relazione per $R$ riportata anche in Haensel et al., eq.~(9.27), \S9.3.
\end{frame}


\subsection{6.10 -- Dinamica orbitale e massa}

\begin{frame}[t]{Perche le binarie permettono di pesare le stelle}
\small
\begin{itemize}
\item L'orbita modula i tempi di arrivo attraverso ritardo geometrico (Roemer), Doppler, dilatazione gravitazionale e relativistica.
\item Dal timing si ricavano periodo orbitale $P_b$, eccentricita $e$, proiezione del semiasse maggiore $x=a_1\sin i/c$ e altri parametri.
\item In teoria newtoniana la \emph{mass function} fornisce una relazione fra $m_1,m_2,i$.
\item Gli effetti post-kepleriani della RG permettono di separare le masse e testare simultaneamente la teoria gravitazionale.
\end{itemize}
\[
M=m_1+m_2,\qquad \nu_b=1/P_b.
\]
\end{frame}

\begin{frame}[t]{Terza legge di Keplero e semiasse maggiore}
\small
Per due masse $m_1,m_2$ in orbita con semiasse maggiore relativo $a=a_1+a_2$:
\[
\left(\frac{2\pi}{P_b}\right)^2
=\frac{G(m_1+m_2)}{a^3}.
\]
Dal baricentro,
\[
m_1a_1=m_2a_2,
\qquad a_1=\frac{m_2}{m_1+m_2}a.
\]
L'orbita vista lungo la linea di vista fornisce dal timing $a_1\sin i$.
\medskip
Il solo moto kepleriano misura combinazioni di masse e inclinazione, non tutte le quantita separatamente.
\end{frame}

\begin{frame}[t]{Derivazione della funzione di massa}
\small
Dalla terza legge:
\[
\frac{4\pi^2a^3}{GP_b^2}=m_1+m_2.
\]
Usando $a_1=m_2a/(m_1+m_2)$ e moltiplicando per $\sin^3 i$:
\[
\frac{4\pi^2(a_1\sin i)^3}{GP_b^2}
=\frac{m_2^3\sin^3 i}{(m_1+m_2)^2}.
\]
Si definisce
\[
f(m_1,m_2,i)=\frac{4\pi^2}{GP_b^2}(a_1\sin i)^3
=\frac{(m_2\sin i)^3}{(m_1+m_2)^2}.
\]
Per $x=a_1\sin i/c$, $f=4\pi^2c^3x^3/(GP_b^2)$.
\end{frame}

\begin{frame}[t]{Vincoli dalla funzione di massa}
\small
\[
f=\frac{m_2^3\sin^3i}{(m_1+m_2)^2}.
\]
\begin{itemize}
\item $f$ e misurata da $P_b$ e $x$.
\item Poiche $\sin i\leq1$, si ottiene un \emph{limite inferiore} a $m_2$ per un dato $m_1$.
\item La conoscenza di $i$ (per esempio da Shapiro) e di un altro parametro post-kepleriano consente spesso di determinare le masse.
\end{itemize}
Se sono osservate entrambe le pulsar:
\[
\frac{m_1}{m_2}=\frac{a_2}{a_1}=\frac{x_2}{x_1}
\]
nel limite dinamico richiesto dal timing relativistico del doppio sistema.
\end{frame}

\begin{frame}[t]{Parametri post-kepleriani: panoramica}
\small
\begin{tabular}{@{}ll@{}}
\toprule
Osservabile & Informazione principale in RG \\
\midrule
$\dot\omega$ (periastron advance) & massa totale $m_1+m_2$ \\
$\gamma_E$ (Einstein delay) & massa, eccentricita e redshift \\
$r,s$ di Shapiro & massa compagna e inclinazione \\
$\dot P_b$ & perdita energetica orbitale (GW) \\
$\Omega_{\rm SO}$ & precessione spin-orbita \\
\bottomrule
\end{tabular}
\medskip
Al crescere del numero di osservabili, il sistema diventa \emph{sovradeterminato}: fissate due masse, le altre misure testano la RG.
\end{frame}

\begin{frame}[t]{Avanzamento del periastron in RG}
\small
Per due masse puntiformi, al primo ordine post-newtoniano:
\[
\dot\omega=3\left(\frac{2\pi}{P_b}\right)^{5/3}
\frac{[G(m_1+m_2)]^{2/3}}{c^2(1-e^2)}.
\]
\begin{itemize}
\item Dipende da massa totale, periodo orbitale ed eccentricita.
\item Misurando $P_b,e,\dot\omega$ si determina $M=m_1+m_2$ nel modello RG.
\item In binarie compatte l'effetto puo essere di molti gradi all'anno, a confronto dei $43''$/secolo di Mercurio.
\end{itemize}
\end{frame}

\begin{frame}[t]{Ritardo di Shapiro: geometria e parametri}
\small
I fotoni attraversano il campo gravitazionale della compagna; il tempo di propagazione aumenta.
Per la parametrizzazione standard del timing di sistemi eccentrici:
\[
\Delta_S=-2r\ln\{1-e\cos u
-s[\sin\omega(\cos u-e)+\sqrt{1-e^2}\cos\omega\sin u]\}.
\]
Con
\[
r=\frac{Gm_2}{c^3},\qquad s=\sin i.
\]
\begin{itemize}
\item $u$ anomalia eccentrica; $\omega$ argomento del periastron.
\item Le orbite viste quasi di taglio ($i\simeq90^\circ$) offrono un segnale Shapiro particolarmente netto.
\end{itemize}
\end{frame}

\begin{frame}[t]{Decadimento orbitale per emissione di GW}
\small
La massa chirp della binaria e
\[
\mathcal M=\frac{(m_1m_2)^{3/5}}{(m_1+m_2)^{1/5}}.
\]
La formula di Peters--Mathews in RG per il decadimento orbitale e
\[
\dot P_b=-\frac{192\pi}{5}
\left(\frac{2\pi G\mathcal M}{c^3P_b}\right)^{5/3}F(e),
\]
con
\[
F(e)=\frac{1+\frac{73}{24}e^2+\frac{37}{96}e^4}
{(1-e^2)^{7/2}}.
\]
La formula riguarda la correzione intrinseca: per il confronto osservativo occorre sottrarre effetti cinematici e galattici.
\end{frame}

\begin{frame}[t]{Figura 6.2 -- Hulse--Taylor e decadimento orbitale}
\small
\begin{figure}[!ht]\includegraphics[trim={0cm 0cm 0 0},clip,keepaspectratio,width=0.95\textwidth]{Schaffner/P-Hulse_Taylor}\label{fig:Schaffner_P-Hulse_Taylor}\end{figure}
\vspace{-2mm}\footnotesize Figura 6.2 del capitolo: spostamento cumulativo dei tempi di periastron dovuto al decadimento orbitale previsto dalla RG (dati storici del testo).
\end{frame}

\begin{frame}[t]{Come $\dot P_b$ verifica l'emissione di onde gravitazionali}
\small
\begin{itemize}
\item In RG il sistema perde energia orbitale: il semiasse maggiore diminuisce e il periodo orbitale si accorcia.
\item Il confronto quantitativo usa masse fissate da altri parametri post-kepleriani.
\item La dipendenza da $e$ e $P_b$ consente un test molto piu stringente di una semplice stima dimensionale.
\item La figura del sistema Hulse--Taylor mostra il ritardo accumulato rispetto a un'orbita a periodo invariabile.
\end{itemize}
\[
\dot P_b<0\quad\Longrightarrow\quad \text{orbita che decade}.
\]
\end{frame}

\subsection{6.11 -- Doppia pulsar e consistenza della RG}

\begin{frame}[t]{Perche J0737--3039 e cosi importante}
\small
Nel sistema doppio sono osservati impulsi di entrambe le stelle di neutroni.
\begin{itemize}
\item Si misurano i parametri orbitali proiettati $x_A,x_B$ di entrambe.
\item Il rapporto di masse e vincolato da $x_B/x_A$ con debolissima dipendenza dal modello gravitazionale ai livelli rilevanti.
\item Una misura post-kepleriana aggiuntiva determina entrambe le masse nell'ipotesi RG.
\item I parametri post-kepleriani residui diventano verifiche indipendenti della coerenza della teoria.
\end{itemize}
\[
\frac{m_A}{m_B}=\frac{x_B}{x_A}.
\]
\end{frame}

\begin{frame}[t]{Figura 6.3 -- Vincoli nel piano delle masse}
\small
\begin{figure}[!ht]\includegraphics[trim={0cm 0cm 0 0},clip,keepaspectratio,width=0.95\textwidth]{Schaffner/doublepulsar_ratio-PKparam}\label{fig:Schaffner_doublepulsar_ratio-PKparam}\end{figure}
\vspace{-2mm}\footnotesize Figura 6.3 del libro: intersezione dei vincoli dinamici e post-kepleriani in una piccola regione del piano $(m_A,m_B)$.
\end{frame}

\begin{frame}[t]{Doppia pulsar: struttura logica del test}
\small
\begin{enumerate}
\item Misurare $P_b,e,x_A,x_B$ dai tempi di arrivo.
\item Ricavare $m_A/m_B=x_B/x_A$.
\item Misurare $\dot\omega$ e determinare $m_A+m_B$ in RG.
\item Confrontare $\gamma_E,r,s,\dot P_b,\Omega_{\rm SO}$ con i valori predetti dalle masse ottenute.
\end{enumerate}
\[
\underbrace{(m_A,m_B)}_{\text{due incognite}}
\quad\stackrel{\text{piu vincoli}}{\longrightarrow}\quad
\text{test della teoria gravitazionale}.
\]
La compatibilita delle linee in un unico punto e la caratteristica fisica della figura.
\end{frame}

\begin{frame}[t]{Momento d'inerzia e sistemi binari relativistici}
\small
Il momento d'inerzia puo influenzare la dinamica orbitale attraverso l'accoppiamento spin--orbita.
\[
J_{\rm spin}=I\Omega_{\rm spin}.
\]
\begin{itemize}
\item Una correzione dell'avanzamento del periastron o della precessione puo contenere il momento angolare di spin della stella.
\item Se la massa e ben determinata, l'inferenza di $I$ vincola direttamente la distribuzione di densita interna e l'EOS.
\item L'effetto e piccolo: occorrono modelli relativistici accurati e timing di precisione.
\end{itemize}
\textbf{Collegamento ai capitoli 4, 7 e 9:} EOS $\to$ soluzione TOV $\to$ $I(M)$ $\to$ osservabili del timing.
\end{frame}

\section{Riepilogo, controlli e limiti delle approssimazioni}

\subsection{Formula sheet e correzioni al capitolo}

\begin{frame}[t]{Relazioni fondamentali: rotazione e campo}
\small
\begin{align*}
\Omega&=\frac{2\pi}P,& \dot\Omega&=-\frac{2\pi\dot P}{P^2},\\
E_{\rm rot}&=\frac12 I\Omega^2,&
\dot E_{\rm sd}&=\frac{4\pi^2I\dot P}{P^3},\\
B_{\rm char,eq}&=3.2\times10^{19}\sqrt{P\dot P}\,\mathrm G,&
\tau_c&=\frac P{2\dot P},\\
n_{\rm br}&=\frac{\Omega\ddot\Omega}{\dot\Omega^2},&
R_L&=\frac{cP}{2\pi}.
\end{align*}
\medskip
$B_{\rm char}$ e $\tau_c$ dipendono dalle ipotesi sul torque: non sono rispettivamente una misura diretta del campo e dell'eta.
\end{frame}

\begin{frame}[t]{Relazioni fondamentali: materia, timing, binarie}
\small
\begin{align*}
I_{\rm Newton}&=\frac{8\pi}{3}\int_0^R\rho r^4dr,&
\mathrm{DM}&=\int_0^Dn_e(s)\,ds,\\
f(m_1,m_2,i)&=\frac{(m_2\sin i)^3}{(m_1+m_2)^2},&
\dot P_b&=-\frac{192\pi}{5}\left(\frac{2\pi G\mathcal M}{c^3P_b}\right)^{5/3}F(e).
\end{align*}
\medskip
Gli sviluppi relativistici del momento d'inerzia e i vincoli post-kepleriani sono indispensabili quando la precisione richiesta supera quella della stima newtoniana.
\end{frame}

\begin{frame}[t]{Correzioni e cautele rispetto alle formule del libro}
\small
\begin{itemize}
\item \textbf{(6.28):} con $\dot\Omega=-K\Omega^n$ il braking index e $n=\Omega\ddot\Omega/\dot\Omega^2$ \emph{senza segno meno}.
\item \textbf{(6.11):} per $\dot P>0$, la \emph{variazione} dell'energia rotazionale ha segno negativo; si usa spesso la potenza positiva $\dot E_{\rm sd}=-\dot E_{\rm rot}$.
\item \textbf{Campo magnetico:} in un dipolo $B_{\rm polo}=2B_{\rm equatore}$; la costante $3.2\times10^{19}$ G adotta il campo equatoriale e il modello di vuoto ortogonale.
\item \textbf{P--$\dot P$:} la death line non e un oggetto universale e dipende dal processo di creazione di coppie.
\item \textbf{Pulsar/GR:} le epoche, i record e i numeri delle tabelle originali hanno valenza storica.
\end{itemize}
\end{frame}

\begin{frame}[t]{Mappa concettuale del capitolo}
\small
\[
\begin{array}{c}
\text{EOS}\ \to\ (M,R,\varepsilon(r),I)\\[2mm]
\downarrow\\[-1mm]
\text{rotazione } (P,\dot P,\ddot P)
\ \to\ (\dot E_{\rm sd},B_{\rm char},\tau_c,n_{\rm br})\\[2mm]
\downarrow\\[-1mm]
\text{magnetosfera e plasma}\ \to\ \text{impulsi radio e death line}\\[2mm]
\downarrow\\[-1mm]
\text{timing binario}\ \to\ \text{masse e test della RG}
\end{array}
\]
\medskip
\textbf{Prospettiva:} il capitolo lega osservabili temporali estremamente accurati a proprieta globali della stella e, indirettamente, alla microfisica della materia densa.
\end{frame}

